\chapter{Outline of design activities} \todo{think about the title}
\label{ch:management}

\todo[inline]{Chapter introduction pending. Remember to connect it to the previous chapters.}


\section{Trade-Off Summary}
\label{sec:trade-off_summary}
\todo[inline]{In this Section, we briefly describe the trade-off performed in the Midterm Report.}

- Arturo


\section{Final Concept Design}
\label{sec:final_concept_design}
\todo[inline]{In this Section, we briefly present the configuration that is going to be designed in more detail in this report.}


\section{Engineering Design and Analysis}
\label{sec:engineering_design_and_analysis}

\todo[inline]{In this Section, we briefly describe the steps required to reach the more detailed design.}

To fulfill its mission as a versatile flight test vehicle, the aircraft is expected to house a range of wing planforms, later referred to as the \textit{planform envelope}. This will allow it to perform a wide range of missions, with example ones listed in \autoref{sec:mission_envelope}. However, the need to house multiple wing geometries is driving the system design. 

Thus, the goals of engineering analysis for the final work package are to design HUGO's fuselage and its associated subsystems, ensuring they accommodate the specified range of planforms, and to develop a standard planform to use when the client does not provide one.

To that end, several analyses described in the Midterm Report \footnote{Midterm Report}: drag estimation, tail sizing, fuel mass fraction computation, and the matching diagram can be re-implemented. For this work package, a structural analysis of both wings and the fuselage is conducted to obtain a complete estimate of aircraft weight. Furthermore, aileron sizing is conducted for each examined planform, and preliminary structural design of the fuselage is performed.

As multiple possible planforms are considered, including multiple possible airfoils, the approach to aerodynamic analysis differs from the one of the Midterm Report \footnote{Midterm Report}. Instead of defining an airfoil and conducting a lifting line analysis to obtain the properties of the three-dimensional wing directly, there is an envelope of possible airfoils defined implicitly by ranges of values of lift slope, maximum lift and moment coefficients. Then, the properties of the three-dimensional wing are deduced from those coefficients using semi-empirical formulae.

Another challenge posed by the planform envelope is that it is a continuous design space spanning several parameters. Evaluating all combinations of those parameters at high resolution is not feasible. As such, a list of parameter combinations considered most constraining is justified in \todo{autoref}.

To allow for general-purpose testing of control systems, the standard planform of HUGO needs to be free of aeroelastic instability. To ensure that, an aeroelastic analysis of the planform is conducted  \todo{autoref}. The standard planform also serves as the benchmark for evaluating the performance of the aircraft in \autoref{ch:performance_analysis}.

The influence of assumptions made during the design process on the ability of HUGO to meet its requirements and on the performance of the system needs to be assessed by means of a sensitivity study. This is addressed in \todo{autoref}.

% 0) The goals for the engineering analysis are to assess whether HUGO can house its specified planforms and to develop the standard planform to be used for control system tests.

% 1) Mention Drag Estimation, the Matching Diagram and the fuel fraction analyses are being reused

% 2) Mention moving away from aerosandbox due to working with airfoils implicitly defined through a range of coefficients

% 3) Mention new 


\section{Boundaries of the Planform Envelope}
\label{sec:Planform envelope boundaries}

The envelope of planforms that HUGO can house needs quantification, both for the user to serve as specifications for wing design, and for further development. The planform researcher has an incentive to test a wing of particular shape, not of particular size (as long as the difference in Reynolds and Froude numbers can be accounted for). As such, they are supposed to select an aspect ratio, as well as taper, dihedral, sweep and thickness-to-chord profiles (those quantities need not be constant along the wing), and wing cross sections along the span, provided the airfoil coefficients of those cross sections fit within the specified range. The wing area and structural thicknesses are deduced from the provided specifications. 

\autoref{tab:planform_envelope} provides the range of allowed planform parameters. A planform is considered compatible with HUGO if it is described by any combination of the parameters within their ranges. To assess whether HUGO can accommodate those possible planforms, the most constraining values of each parameter are identified, and planforms represented by combinations of those critical values are analysed for compliance with requirements. Overall, there are $4\cdot2\cdot2\cdot3\cdot2 = 96$ planforms considered the critical points of the planform envelope. As for the justification of the parameter ranges and critical values:
\begin{itemize}
    \item $AR$: the lower bound corresponds to low-$AR$ aerobatic aircraft \footnote{\label{foot:Aerobatic_AR}\href{https://web.archive.org/web/20060414172013/http://www.sukhoi.org/eng/planes/civil/su-26/characteristic/}{Sukhoi Su-26 Specifications}}, while the upper bound is as considered in previous reports \footnote{Baseline}\footnote{Midterm}. As increasing the $AR$ improves the aerodynamic efficiency while increasing wing weight, a single value of this parameter cannot be considered constraining. Therefore, four values from across the range are included as critical.
    \item $\lambda$: any taper profile that results in the wing not widening while moving inboard, but preventing it from becoming excessively thin, is considered acceptable. In the considered range, the most constraining case, when both aerodynamics and structural performance are negatively impacted, is when $\lambda$ is 1 \cite[p.401-405, p.438]{raymer1992aircraft, AndersonAerodynamics}.
    \item $\Lambda_{c/4}$: positive sweeps up to slightly above the Midterm Report value \footnote{Midterm Report}. The extremum values of sweep are taken as critical, as for maximum sweep the effective airspeed is reduced, and the aerodynamic centre is most aft, and for the minimum sweep the aerodynamic centre is most fore. The former is challenging for tail sizing for stability and main wing performance, while for the latter, empennage sizing for controllability is compromised. 
    \item $t/c$: The range encompasses thin and potentially unstable wings and high-thickness profiles similar to those at wing roots of large aircraft \footnote{\label{foot:reference_tc_values}\href{https://m-selig.ae.illinois.edu/ads/aircraft.html}{David Lednicer, The Incomplete Guide to Airfoil Usage}}. As for the critical points, parasitic drag increases with $t/c$ \cite[p.283]{raymer1992aircraft}, but the wing weight decreases significantly \cite[p.401-405]{raymer1992aircraft}. Thus, three values across the parameter range are selected for evaluation.
    \item $C_{l_\text{max}}$: To ensure sufficient low-speed performance, a minimum value of this parameter is established as 1, which is below the $C_{l_\text{max}}$ of most aifoils \footnote{\label{foot:NACA_4415}\href{http://airfoiltools.com/airfoil/details?airfoil=naca4415-il}{Airfoil Tools, NACA 4415}}\footnote{\label{foot:NACA_24112}\href{http://airfoiltools.com/airfoil/details?airfoil=naca24112-jf}{Airfoil Tools, NACA 24112}}.
    \item $C_{m_{ac}}$: The extrema values of this parameter are most constraining for sizing the empennage for controllability and, as such, taken as critical points. The minimum $C_{m_{ac}}$ accounts for high camber sections such as NACA 4415 with $C_{m_{ac}}$ of around -0.12 \footref{foot:NACA_4415}, while the upper bound accounts for reflexed airfoils such as NACA 24112  with $C_{m_{ac}}$ of up to 0.03 \footref{foot:NACA_24112}.
    \item $C_{l_0}$: The impact of this parameter is secondary - provided $C_{l_\alpha}$ and $C_{l_{\mathrm{max}}}$ are satisfactory, it only changes the angle of attack required for a given amount of lift. Hence, any non-negative vlue of the parameter is allowed, with zero taken as the critical value.
    \item $C_{l_\alpha}$: The lift slope for thin airfoil is typically close to $2\pi$ and this value is taken as critical.
    \item Empennage type: both replaceable tail cone and canard port are available at the aircraft, and they may be used simultaneously. As the moment arms of both the tail and the canard are fixed, cases when only the canard or only the tail are present are taken as constraining.
    \item $\Gamma$: The lower bound of dihedral at zero is set not to worsen the wing tip clearance, and the upper bound is where the lift reduction from dihedral becomes approximately 1\%. Due to its secondary influence, $\Gamma$ is excluded from the analysis, with a value of 0 used for landing gear clearance computations.
\end{itemize}

Not every combination of wing cross section and sweep is expected to reach Mach 0.75. However, provided that shocks do not develop up to that speed, a planform whose cross section coefficients are within the stated range is supposed to be able to reach Mach 0.75. Also, there is a possibility that a planform whose parameters are outside of the stated ranges is compatible with HUGO.

\begin{table}[]
    \centering
    \caption{Overview of the Planform envelope. For each parameter, the full range of it considered compatible with HUGO is provided, along with critical values that the requirements need to be assessed for. $f(y)$ means the parameter can vary spanswise. }
    \begin{tabular}{|c|c|c||c|c|c|}\hline
         \textbf{Parmeter}& \textbf{Range} & \textbf{Critical Values}&\textbf{Parmeter}& \textbf{Range} & \textbf{Critical Values}  \\\hline
         $AR$ & [5; 27]  & 5, 10, 17, 27 & $\lambda$ & $f(y)\in[0.1; 1]$ & 1 \\\hline
         $\Lambda_{c/4}$ & $f(y)\in[0; 20^\circ]$  & 0, 20$^\circ$ & $t/c$ & $f(y)\in[0.08,\ 0.18]$ & 0.08, 0.12, 0.18 \\\hline
         $C_{l_{\textbf{max}}}$ & $f(y)\in[1, \infty)$  & 1 & $C_{m_{ac}}$ & $f(y)\in[-0.15,\ 0.1]$ & -0.15, 0.1 \\\hline
         $C_{l_{0}}$ & $f(y)\in[0,\ \infty)$  & 0 & $C_{l_{\alpha}}$ & $f(y)\in[2\pi,\ \infty)$ & 2$\pi$ \\\hline
         Empennage types & Canard and/or tail  & Canard or tail & $\Gamma$ & $f(y)\in[0,\  8^\circ]$ & 0 \\\hline
    \end{tabular}
    
    \label{tab:planform_envelope}
\end{table}

\section{The standard planform}
\begin{table}[h]
    \centering
    \caption{The geometry of the standard planform}
    \begin{tabular}{|c|c||c|c||c|c||c|c|}\hline
         \textbf{Parameter}&\textbf{Value}&\textbf{Parameter}&\textbf{Value}&\textbf{Parameter}&\textbf{Value}&\textbf{Parameter}&\textbf{Value}  \\\hline
         $AR$& 27 & $\lambda$ & 0.3 & $\Lambda_{c/4}$ & 15$^\circ$ & $\Gamma$ & 0\\\hline
         $b$ &3.2 m & $t/c$ & 0.12 & airfoil & NASA SC(2)-0012 & Empennage & T-tail\\\hline
    \end{tabular}
    
    \label{tab:standard_wing_planform}
\end{table}

The geometry of the standard planform, to be used when the client does not provide one of their own, is presented in \autoref{tab:standard_wing_planform}. The planform is based on the one obtained in the Midterm Report \footnote{Midterm Report}. It is adapted for slightly larger aircraft mass and aeroelastic stability\todo{autoref}, which required adjusting the span and the taper ratio, respectively.



\todo[inline]{Chapter conclusion pending. Remember to connect it to the chapters that come after.}